% Part 5 – Theory of Score–Based Generative Models

% ---------------------------------------------------------------------------
% This section introduces the theoretical underpinnings of score–based
% generative models.  We build upon the notion of the score function
% introduced in the energy–based model context and show how it can be
% estimated without computing the normalising constant.  We discuss
% score matching, denoising score matching, and multi–scale noise
% perturbation.  We also explain why the divergence term in the original
% score matching objective is problematic in high dimensions and how
% multi–scale perturbations circumvent this issue.  Connections to
% stochastic differential equations and diffusion processes are sketched to
% motivate the sampling algorithms used later in the implementation.

\section{Score--Based Models: Theory}

\subsection{From Energies to Scores}

In Part~2 we defined the score function $\mathbf{s}_\theta(\mathbf{x}) =
\nabla_{\mathbf{x}} \log p_\theta(\mathbf{x})$ and observed that it can be
expressed as $-\nabla_{\mathbf{x}} E_\theta(\mathbf{x})$ because the gradient
of the partition function is zero.  For
score–based models we abandon the energy view entirely and instead seek to
learn the score function directly.  This perspective has several
advantages:

\begin{itemize}[leftmargin=*]
  \item The score function depends only on the gradient of the log–density
    with respect to $\mathbf{x}$ and is independent of the normalising
    constant $Z_\theta$, obviating the need to evaluate intractable
    integrals.
  \item Once a good estimate of the score function $\mathbf{s}_\theta$ is
    available, one can perform sampling using Langevin dynamics on the
    estimated gradient field.  The stationary distribution of the resulting
    Langevin SDE is the target density (up to approximation error), even
    though the density itself is not explicitly known.
  \item Learning the score function is inherently flexible: one can model
    gradients with arbitrarily expressive neural networks without imposing
    global normalisation constraints.
\end{itemize}

Given a target density $p_{\text{data}}(\mathbf{x})$ (unknown up to the
data samples), the goal is to learn a parameterised score function
$\mathbf{s}_\theta(\mathbf{x})$ that approximates $\nabla_{\mathbf{x}} \log
p_{\text{data}}(\mathbf{x})$ for most $\mathbf{x}$ in the support of the
data distribution.  Several estimators have been proposed; we focus on
score matching and its variants.

\subsection{Score Matching and Its Limitations}

\emph{Score matching}, introduced by Hyv\"arinen, provides a way to
estimate the score function by minimising the Fisher divergence between the
model score and the true data score:
\begin{equation}
    \mathcal{J}(\theta) = \frac{1}{2} \int p_{\text{data}}(\mathbf{x})
    \big\| \mathbf{s}_\theta(\mathbf{x}) - \nabla_{\mathbf{x}}
    \log p_{\text{data}}(\mathbf{x}) \big\|_2^2 \,\mathrm{d}\mathbf{x}.
    \label{eq:fisher-divergence-sc}
\end{equation}
Differentiating Eq.~\eqref{eq:fisher-divergence-sc} with respect to $\theta$
and integrating by parts yields an equivalent loss that does not
explicitly depend on the unknown data score but does involve the divergence
of the model score:
\begin{equation}
    \mathcal{J}(\theta) = \mathbb{E}_{\mathbf{x}\sim p_{\text{data}}}
    \big[ \tfrac{1}{2} \| \mathbf{s}_\theta(\mathbf{x}) \|^2
    + \nabla\cdot\mathbf{s}_\theta(\mathbf{x}) \big].
    \label{eq:score-matching-loss-sc}
\end{equation}
While Eq.~\eqref{eq:score-matching-loss-sc} removes the unknown data score,
it introduces the divergence term $\nabla\cdot\mathbf{s}_\theta(\mathbf{x}) =
\sum_i \partial s_i(\mathbf{x}) / \partial x_i$, which requires computing
second derivatives of the model output with respect to its input.  For
high–dimensional data (e.g., images with thousands of pixels) this
computation is expensive and often unstable.  Moreover, it imposes
stringent smoothness requirements on the neural network architecture.  As a
result, direct score matching has seen limited practical adoption in deep
learning.

\subsection{Denoising Score Matching}

To avoid the divergence term and the dependence on the unknown data
score, Vincent proposed \emph{denoising score matching} (DSM).  The key
idea is to corrupt each data point with Gaussian noise and train a
network to predict the denoising direction.  Let $\tilde{\mathbf{x}}$ be
obtained from a clean sample $\mathbf{x}$ via
\begin{equation}
    \tilde{\mathbf{x}} = \mathbf{x} + \sigma \boldsymbol{\epsilon},
    \quad \boldsymbol{\epsilon} \sim \mathcal{N}(\mathbf{0},\mathbf{I}),
    \label{eq:dsm-perturb}
\end{equation}
where $\sigma>0$ is a user–chosen noise level.  The distribution
of $\tilde{\mathbf{x}}$ is the convolution of the data distribution with
a Gaussian kernel.  A remarkable result shows that the score of this noisy
distribution is proportional to the negative of the injected noise,
namely
\begin{equation}
    \nabla_{\tilde{\mathbf{x}}} \log p_\sigma(\tilde{\mathbf{x}}) =
    -\frac{\boldsymbol{\epsilon}}{\sigma}.
\end{equation}
DSM therefore trains a parameterised score function
$\mathbf{s}_\theta(\tilde{\mathbf{x}},\sigma)$ to predict $-\boldsymbol{\epsilon}/\sigma$.
The loss function for DSM can be written as
\begin{equation}
    \mathcal{L}_{\text{DSM}}(\theta) =
    \mathbb{E}_{\mathbf{x},\boldsymbol{\epsilon}}
    \Big[ \tfrac{1}{2} \sigma^2 \big\| \mathbf{s}_\theta(\mathbf{x} +
    \sigma \boldsymbol{\epsilon}, \sigma) + \tfrac{\boldsymbol{\epsilon}}{\sigma}
    \big\|^2 \Big].
    \label{eq:dsm-loss-sec}
\end{equation}
Because the target $-\boldsymbol{\epsilon}/\sigma$ is known, DSM turns
score estimation into a supervised learning problem.  Crucially, the
loss in Eq.~\eqref{eq:dsm-loss-sec} does not involve the divergence
operator and avoids computing second derivatives.  As $\sigma$ tends to
zero, the DSM objective approaches the original score matching loss, but
training at zero noise is unstable and provides no gradient signal away
from the data manifold.  Adding a small amount of noise spreads the
support of the data distribution and yields well–behaved gradients.

\subsection{Multi--Scale Noise Perturbation and Weighted DSM}

Using a single noise level $\sigma$ allows the model to learn the score
of the noisy density $p_\sigma$, but it does not directly provide the
score of the clean distribution.  To bridge this gap, Song and Ermon
introduced a \emph{multi–scale noise schedule}.  They select a sequence
of noise levels $\sigma_1 > \sigma_2 > \cdots > \sigma_L$ that spans a
large range, typically on a logarithmic scale.  During training, each
mini–batch is corrupted with a noise level sampled uniformly from this
set, and the noise level itself is provided as an input to the network.

To ensure that contributions from different noise scales are balanced,
the DSM loss is weighted by $\sigma^2$:
\begin{equation}
    \mathcal{L}_{\text{MDSM}}(\theta) =
    \mathbb{E}_{\mathbf{x},\boldsymbol{\epsilon}}
    \mathbb{E}_{\sigma \sim \mathcal{U}(\{\sigma_1,\dots,\sigma_L\})}
    \Big[ \tfrac{1}{2} \sigma^2 \big\|\mathbf{s}_\theta(\mathbf{x} +
    \sigma \boldsymbol{\epsilon}, \sigma) + \tfrac{\boldsymbol{\epsilon}}{\sigma}
    \big\|^2 \Big].
    \label{eq:mdsm-loss}
\end{equation}
Here $\mathcal{U}(\{\sigma_1,\dots,\sigma_L\})$ denotes the uniform
distribution over the discrete ladder.  Large noise levels dominate the
loss due to the $\sigma^2$ weighting, forcing the model to focus on
coarse–scale structures.  As the noise decreases, the model learns finer
details.  This multi–scale denoising score matching (MDSM) objective
provides gradients across the entire range of resolutions and is critical
for high–quality generation.
\subsection{Annealed Langevin Dynamics and Connections to Diffusion Processes}

Once the score network has been trained across multiple noise levels,
sampling proceeds by reversing the corruption process.  We start from a
Gaussian distribution with variance $\sigma_L^2$ (the largest noise
level) and iteratively refine the sample via \emph{annealed Langevin
dynamics}.  For each level $\sigma_k$, we perform a fixed number of
Langevin updates
\begin{equation}
    \mathbf{x}_{t+1} = \mathbf{x}_t + \eta_k\,\mathbf{s}_\theta(\mathbf{x}_t,
    \sigma_k) + \sqrt{2\eta_k}\,\boldsymbol{\xi}_t,
\end{equation}
where $\eta_k$ is a step size proportional to $(\sigma_k/\sigma_1)^2$ and
$\boldsymbol{\xi}_t \sim \mathcal{N}(\mathbf{0},\mathbf{I})$.  After $T$
steps at level $\sigma_k$ the noise level is reduced to $\sigma_{k-1}$
and the process repeats.  In the limit of infinitely many noise scales
and infinitesimal step sizes, this discrete procedure approximates the
solution of a reverse stochastic differential equation.  Song,
Sohl–Dickstein and collaborators formalised this connection and showed
that score–based models can be interpreted as simulating a diffusion
process backwards in time using the learned score function.  Although the
continuous SDE formulation provides a clear theoretical framework, the
discrete ALD algorithm is simple to implement and performs well in
practice.  We adopt this algorithm in our implementation.

\subsection{Mixture Model Revisited and Mode Trapping}

The mixture–of–Gaussians example from Part~2 illustrates a limitation of
gradient–based samplers, whether they operate on energies or scores.
Consider again the mixture density
\begin{equation}
    p(\mathbf{x}) = w_a p_a(\mathbf{x}) + w_b p_b(\mathbf{x}),
\end{equation}
where the supports of $p_a$ and $p_b$ are disjoint.  The score within
each region depends only on the local component and ignores the mixing
weights.  Consequently, if we initialise Langevin dynamics or ALD within
component $a$, the trajectory will remain in the domain of $p_a$ because
the gradient of the log–density vanishes at the boundary.  The sampler
cannot cross the energy barrier separating the modes.  This phenomenon is
often referred to as \emph{mode trapping}.  In practice, the multi–scale
noise schedule helps to some extent: at large noise levels the Gaussian
convolution blends the components and the score field becomes smoother,
allowing the chain to cross between modes.  However, components with very
small mixing weights may still be under–sampled.  Designing samplers that
explore all modes proportionally remains an open research problem.

\subsection{Summary of Score–Based Theory}

This part has developed the theoretical foundations of score–based
generative modeling.  Starting from the score function we derived the
Fisher divergence loss and explained why its naive implementation is
impractical for high–dimensional data.  We introduced denoising score
matching, which circumvents the need for divergence terms by adding
Gaussian noise to the data and casting score estimation as a supervised
learning problem.  A multi–scale noise ladder and corresponding
weighting scheme were presented to train the model across a range of
resolutions.  We then described annealed Langevin dynamics, a discrete
approximation to the reverse SDE used for sampling, and discussed its
connection to diffusion processes.  Finally, we revisited the
mixture–of–Gaussians example to highlight the mode trapping issue and
underscored the need for careful sampler design.  These insights lay the
groundwork for the practical implementation of noise conditional score
networks in the next part.
